\documentclass[a4paper,11pt]{article}
\usepackage{exam-style-341}
\examinehead{341 农业综合知识三 · 网络技术与应用}{模拟试卷（二）参考答案}
\begin{document}

\answercover{341 农业综合知识三 · 网络技术与应用}{模拟试卷（二）参考答案}

\answerpart{一、单项选择题}
\setcounter{question}{0}
\q B。\quad \q B。\quad \q B。\quad \q A。\quad \q B。\quad \q C。\quad \q B。\quad \q B。\quad \q C。\quad \q B。
\begin{analysis}
第 1 题：计算机网络的主要功能是数据通信和资源共享（硬件、软件、数据资源），「完全替代邮政通信系统」「只能进行分布式计算」都过于绝对，故 B 正确。
第 2 题：总线型拓扑用一条公共传输媒体连接所有结点，结构简单、成本低、易于扩充；其缺点是可靠性较差、故障定位困难、不易实现流量控制，故 A、C、D 错误。
第 3 题：TCP 是面向字节流的，序号和确认号都以字节为单位。确认号表示「期望收到对方下一个报文段的第一个数据字节的序号」，即 \(500+200=700\)，故选 B（不是 701，也不等于序号 500）。
第 4 题：IP 地址标识主机（接口）与所在网络的连接，主机移动到另一网络时 IP 地址必须改变，而 MAC 地址由网卡硬件决定、保持不变；路由器转发时只改写 MAC 帧的地址，不改变分组的源、目的 IP 地址；路由器的每个接口都必须有 IP 地址，故只有 A 正确。
第 5 题：分组交换把长报文划分为较小的分组并采用存储转发，分组可经不同路径并行传输，因而线路利用率高、时延较小；但分组要附加首部，会产生额外开销，B 正确。
第 6 题：TCP 的发送窗口 \(=\) 接收窗口与拥塞窗口中的较小者（\(W=\min(rwnd,cwnd)\)），既受接收方接收能力限制，也受网络拥塞程度限制，故选 C。
第 7 题：FTP、HTTP、SMTP、Telnet 都基于 TCP；DNS 主要用于 UDP 的 53 端口，仅在区域传送等少数场合使用 TCP，故 B 正确。
第 8 题：\(/22\) 的前缀长度为 22 位，主机号位数 \(=32-22=10\)，地址块内共有 \(2^{10}=1024\) 个地址，除去主机号全 0 的网络地址和主机号全 1 的广播地址，可分配给主机的地址为 \(2^{10}-2=1022\) 个。注意：「/22 包含 4 个 C 类网络」说的是 1024 个**地址**，可分配主机数要再减 2。
第 9 题：交换机工作于数据链路层，采用存储转发方式，收到完整帧后根据帧首部中的目的 MAC 地址查交换表（转发表）并从相应端口转发；帧的内容不会改变；交换机用于交换式以太网而非共享式总线；只有目的地址未知或为广播地址时才向所有端口（除输入端口外）转发，故 C 正确。
第 10 题：最短帧长 \(=\) 争用期 \(\times\) 数据率 \(=51.2\times 10^{-6}\ \mathrm{s}\times 10\times 10^{6}\ \mathrm{bit/s}=512\ \mathrm{bit}=512/8=64\) 字节，故选 B。46 字节是数据字段的最小长度（\(64-18\)），不要混淆。
\end{analysis}

\answerpart{二、填空题}
\setcounter{question}{0}
\q 同步（时序）；语法。【两个空顺序可互换，写「语法；同步（时序）」同样给分。】\quad
\q 网际层（IP 层）；应用层。【两个空顺序可互换。】\quad
\q 确认（ACK）；重传（超时重传）。\quad
\q 比特（二进制位）；比特率 \(=\) 波特率 \(\times\) 一个码元携带的比特数（即 \(C=B\log_2 n\)）。【第一空填「比特（二进制位）」；第二空写出关系式即可，也可写 \(C=B\log_2 n\)、\(B=C/\log_2 n\)。】\quad
\q 快重传；快恢复。
\begin{analysis}
第 1 题：语法规定数据与控制信息的格式；语义规定发出何种控制信息、完成何种动作以及做出何种响应；同步（时序）规定事件实现顺序的详细说明。
第 2 题：TCP/IP 的四层由下到上为网络接口层（网络接口层/链路层）、网际层（IP 层）、运输层、应用层。
第 3 题：停止等待协议中，发送方每发送一个数据帧就停止发送、等待确认；超时未收到确认就重传该帧，用序号区分新帧与重传帧，否则无法判断收到的是新帧还是重传帧的副本。
第 4 题：比特率是每秒传输的比特数；若一个码元携带 \(n\) 个比特，则比特率 \(C\) 与波特率 \(B\) 的关系为 \(C=B\log_2 n\)（当 \(n=1\) 时二者相等，故不能说比特率一定等于波特率）。
第 5 题：TCP 拥塞控制的四种算法是慢开始、拥塞避免、快重传和快恢复；其中快重传要求接收方收到失序报文段时立即发出重复确认，快恢复则把慢开始门限减半后直接进入拥塞避免。
\end{analysis}

\answerpart{三、名词解释}
\setcounter{question}{0}
\q \textbf{服务访问点（SAP，Service Access Point）}：同一系统中相邻两层实体进行交互（交换信息）的地方，是上层使用下层所提供服务的逻辑接口，实际上就是下层向上层提供服务的入口。每个 SAP 都有一个能唯一标识它的地址（如运输层的 SAP 就是端口号）；一个下层实体可以同时为多个上层实体服务，也就是在同一层中存在多个 SAP。协议是「水平方向」对等层之间的通信规则，而服务是通过 SAP 在「垂直方向」由下层向上层提供的。

\q \textbf{MAC 地址（物理地址、硬件地址）}：数据链路层使用的地址，又称物理地址或硬件地址，用 48 位（6 字节）二进制数表示，通常写成 12 个十六进制数字，如 \texttt{00-1A-2B-3C-4D-5E}。它固化在网卡（网络适配器）的 ROM 中，由 IEEE 统一分配，前 3 个字节为组织唯一标识符（OUI）、后 3 个字节由厂商分配，因此每个网卡的 MAC 地址在全世界范围内是唯一的（可视为固定的、不随主机所在网络改变）。以太网用 MAC 地址标识同一局域网上的每个接口，帧的接收与转发都以 MAC 地址为依据。

\q \textbf{网络协议}：为进行网络中的数据交换而建立的规则、标准或约定，是计算机网络中不可缺少的组成部分。网络协议由三个要素组成：语法（数据与控制信息的格式）、语义（需要发出何种控制信息、完成何种动作以及做出何种响应）、同步（事件实现顺序的详细说明）。协议是分层的，每一层都有自己的协议，对等层之间按协议交换数据单元；协议是「水平的」，即协议是两台（或多台）对等实体之间通信规则的集合。

\q \textbf{运输层的端口}：运输层为标识一台主机中参与通信的各个应用进程而使用的协议端口号，简称端口。端口用一个 16 位二进制数（\(0\sim65535\)）表示，它是应用进程与运输层实体交互的服务访问点；运输层用「IP 地址 \(+\) 端口号」共同确定互联网中的一个进程（套接字 Socket）。端口分为三类：熟知端口（\(0\sim1023\)，如 FTP 的 21、HTTP 的 80、DNS 的 53）、登记端口（\(1024\sim49151\)）和客户端使用的短暂端口（\(49152\sim65535\)）。

\answerpart{四、简答与计算题}
\setcounter{question}{0}

\q \textbf{解：}

\textbf{（1）按层次划分的域名服务器（3 分）}
\begin{enumerate}[leftmargin=2em, itemsep=2pt, topsep=3pt]
  \item \textbf{根域名服务器}：最高层次的域名服务器，知道所有顶级域名服务器的域名和 IP 地址。全世界共有 13 个根域名服务器（用 A\(\sim\)M 编号，其中不少是冗余分布的集群）；本地域名服务器无法解析域名时，首先向根域名服务器查询。
  \item \textbf{顶级域名服务器（TLD）}：负责管理在该顶级域名（如 com、cn、edu）下注册的所有二级域名，收到查询时给出相应权限域名服务器（或下一级域名服务器）的地址。
  \item \textbf{权限域名服务器（授权域名服务器）}：负责管理一个「区」（zone）的域名服务器，保存该区内所有主机的域名到 IP 地址的映射，能够把管辖范围内的域名直接解析为 IP 地址，因此它是最终给出权威答案的服务器。
  \item \textbf{本地域名服务器（默认域名服务器）}：不属于上述层次结构，但非常重要。当一个主机发出 DNS 查询请求时，这个查询报文首先发送给本地域名服务器；本地域名服务器通常离用户较近、缓存了大量最近查询过的域名映射，是主机进行域名解析的「代理人」。
\end{enumerate}

\textbf{（2）解析过程中的关系（2 分）}
主机与本地域名服务器之间采用\textbf{递归查询}：主机只要发出一次查询，本地域名服务器就必须代为查出结果（或返回查询失败）。
本地域名服务器与根域名服务器、顶级域名服务器、权限域名服务器之间通常采用\textbf{迭代查询}：根域名服务器收到查询后不代为继续查询，而是返回「下一步应向哪个顶级域名服务器的 IP 地址查询」；本地域名服务器再向该顶级域名服务器查询，如此逐级向下，直到某个权限域名服务器给出所要的 IP 地址（或返回查询失败）。此外，凡是有域名服务器收到查询，都会把结果保存在高速缓存中并设置生存时间（TTL），以便后续查询直接回答，从而减少查询次数、提高效率。

\begin{analysis}
本题评分要点：①正确列出四类服务器——根域名服务器、顶级域名服务器、权限（授权）域名服务器、本地域名服务器（每类约 0.75 分，共 3 分）；②说明主机到本地域名服务器为递归查询、本地域名服务器到根/顶级/权限域名服务器通常为迭代查询，并提到高速缓存的作用（2 分）。
\end{analysis}

\q \textbf{解：}

\textbf{（1）每个分片最多携带的数据字节数（1.5 分）}

IP 数据报首部长 20 字节，链路 MTU（最大传送单元）限定「首部 \(+\) 数据」的总长度不超过 1420 字节，因此每个分片可携带的数据最多为：
\[
1420-20=1400\ \text{字节}
\]
但由于 IP 首部中的\textbf{片偏移字段以 8 字节为基本单位}（片偏移 \(=\) 本片数据在原始数据报数据部分中的起始位置 \(\div 8\)），每个分片携带的数据长度必须是 8 的整数倍，而 \(1400\div 8=175\) 恰好整除，所以每个分片最多可携带 1400 字节数据。若有余数（如 MTU 恰使可用长度为 1401 字节），则必须向下取 8 的最大整数倍。

\textbf{（2）分片个数与片偏移（2.5 分）}

原始数据报首部 20 字节，故其数据部分长度为：
\[
4000-20=3980\ \text{字节}
\]
按每片最多携带 1400 字节数据划分：
\[
3980=1400+1400+1180
\]
因此需要划分为 \(\mathbf{3}\) 个分片。片偏移字段的值 \(=\) 该片第一个数据字节在原始数据部分中的序号 \(\div 8\)：
\begin{center}
\begin{tabular}{ccccc}
\toprule
分片 & 携带数据字节数 & 数据起始位置（字节序号） & 片偏移 \(=\) 起始位置/8 & 总长度（字节） \\
\midrule
第 1 片 & 1400 & 0     & \(0/8=0\)        & \(1400+20=1420\) \\
第 2 片 & 1400 & 1400  & \(1400/8=175\)   & \(1400+20=1420\) \\
第 3 片 & 1180 & 2800  & \(2800/8=350\)   & \(1180+20=1200\) \\
\bottomrule
\end{tabular}
\end{center}
故三个分片的片偏移字段取值分别为 \(\mathbf{0}\)、\(\mathbf{175}\)、\(\mathbf{350}\)。

\textbf{（3）MF 标志与总长度（1 分）}
\begin{itemize}[leftmargin=2em, itemsep=2pt, topsep=3pt]
  \item 第 1 片：MF \(=1\)（后面还有分片），总长度 \(=1420\) 字节；
  \item 第 2 片：MF \(=1\)（后面还有分片），总长度 \(=1420\) 字节；
  \item 第 3 片：MF \(=0\)（最后一片），总长度 \(=1200\) 字节。
\end{itemize}
各分片的 DF（不分片）标志均为 0（本题未禁止分片），且各分片在通过 MTU 为 1420 字节的链路时，总长度均不超过 1420 字节。

\textbf{验算}：\(1400+1400+1180=3980\) 字节，与原始数据报数据部分长度相等；片偏移 \(350\times 8=2800\)，\(2800+1180=3980\)，说明最后一片正好覆盖到原始数据部分的末尾；每个分片的总长度都不超过 MTU \(=1420\) 字节。

\begin{analysis}
本题评分要点：①说明分片最多携带 \(1420-20=1400\) 字节，并明确指出「片偏移以 8 字节为单位，故每片数据长度必须是 8 的整数倍」（1.5 分）；②由 \(4000-20=3980=1400+1400+1180\) 得出 3 个分片（1 分），片偏移分别为 \(0/8=0\)、\(1400/8=175\)、\(2800/8=350\)（1.5 分）；③正确写出 MF 依次为 \(1,1,0\)，总长度依次为 \(1420,1420,1200\) 字节（1 分）。常见错误：把片偏移写成字节序号 1400、2800；把总长度写成 4000 或忘记加 20 字节首部；忽略「片偏移以 8 字节为单位」而使用非 8 整数倍的分片数据长度。
\end{analysis}

\end{document}
